\documentclass[10pt,a4paper]{amsart}
\usepackage{amsmath,amssymb,amsthm,graphicx}
% For commutative diagrams we use Xy-pic
\usepackage[all]{xy}

% We use 2cell for 2-commutative diagrams.
\xyoption{2cell}
% Theorem environments.
%
\theoremstyle{plain}
\newtheorem{theorem}[subsection]{Theorem}
\newtheorem{proposition}[subsection]{Proposition}
\newtheorem{lemma}[subsection]{Lemma}

\theoremstyle{definition}
\newtheorem{definition}[subsection]{Definition}
\newtheorem{example}[subsection]{Example}
\newtheorem{exercise}[subsection]{Exercise}
\newtheorem{situation}[subsection]{Situation}

\theoremstyle{remark}
\newtheorem{remark}[subsection]{Remark}
\newtheorem{remarks}[subsection]{Remarks}

\numberwithin{equation}{subsection}

% Macros
%
\def\lim{\mathop{\mathrm{lim}}\nolimits}
\def\colim{\mathop{\mathrm{colim}}\nolimits}
\def\Spec{\mathop{\mathrm{Spec}}}
\def\Hom{\mathop{\mathrm{Hom}}\nolimits}
\def\Ext{\mathop{\mathrm{Ext}}\nolimits}
\def\SheafHom{\mathop{\mathcal{H}\!\mathit{om}}\nolimits}
\def\SheafExt{\mathop{\mathcal{E}\!\mathit{xt}}\nolimits}
\def\Sch{\mathit{Sch}}
\def\Mor{\mathop{\mathrm{Mor}}\nolimits}
\def\Ob{\mathop{\mathrm{Ob}}\nolimits}
\def\Sh{\mathop{\mathit{Sh}}\nolimits}
\def\NL{\mathop{N\!L}\nolimits}
\def\CH{\mathop{\mathrm{CH}}\nolimits}
\def\proetale{{pro\text{-}\acute{e}tale}}
\def\etale{{\acute{e}tale}}
\def\QCoh{\mathit{QCoh}}
\def\Ker{\mathop{\mathrm{Ker}}}
\def\Im{\mathop{\mathrm{Im}}}
\def\Coker{\mathop{\mathrm{Coker}}}
\def\Coim{\mathop{\mathrm{Coim}}}

% Boxtimes
%
\DeclareMathSymbol{\boxtimes}{\mathbin}{AMSa}{"02}

%
% Macros for moduli stacks/spaces
%
\def\QCohstack{\mathcal{QC}\!\mathit{oh}}
\def\Cohstack{\mathcal{C}\!\mathit{oh}}
\def\Spacesstack{\mathcal{S}\!\mathit{paces}}
\def\Quotfunctor{\mathrm{Quot}}
\def\Hilbfunctor{\mathrm{Hilb}}
\def\Curvesstack{\mathcal{C}\!\mathit{urves}}
\def\Polarizedstack{\mathcal{P}\!\mathit{olarized}}
\def\Complexesstack{\mathcal{C}\!\mathit{omplexes}}
% \Pic is the operator that assigns to X its picard group, usage \Pic(X)
% \Picardstack_{X/B} denotes the Picard stack of X over B
% \Picardfunctor_{X/B} denotes the Picard functor of X over B
\def\Pic{\mathop{\mathrm{Pic}}\nolimits}
\def\Picardstack{\mathcal{P}\!\mathit{ic}}
\def\Picardfunctor{\mathrm{Pic}}
\def\Deformationcategory{\mathcal{D}\!\mathit{ef}}
\usepackage{cleveref}
\crefname{originallemma}{original lemma}{original lemmas}
\Crefname{originallemma}{Original lemma}{Original lemmas}
\makeatletter
\newlabel{limits-lemma-integral-limit-finite-and-finite-presentation}{{09YZ}{source}{}{sourcepointer-1}{}}
\newlabel{limits-lemma-integral-limit-finite-and-finite-presentation@cref}{{[originallemma][0][]09YZ}{[1][0][]source}{}{}{}}
\makeatother
\crefname{originallemma}{original lemma}{original lemmas}
\Crefname{originallemma}{Original lemma}{Original lemmas}
\makeatletter
\newlabel{lemma-linus-hamann}{{0EYN}{source}{}{sourcepointer-2}{}}
\newlabel{lemma-linus-hamann@cref}{{[originallemma][0][]0EYN}{[1][0][]source}{}{}{}}
\makeatother
\crefname{originallemma}{original lemma}{original lemmas}
\Crefname{originallemma}{Original lemma}{Original lemmas}
\makeatletter
\newlabel{lemma-relative-colimit}{{09Z1}{source}{}{sourcepointer-3}{}}
\newlabel{lemma-relative-colimit@cref}{{[originallemma][0][]09Z1}{[1][0][]source}{}{}{}}
\makeatother
\crefname{originallemma}{original lemma}{original lemmas}
\Crefname{originallemma}{Original lemma}{Original lemmas}
\makeatletter
\newlabel{lemma-finite-pushforward-commutes-with-base-change}{{0959}{source}{}{sourcepointer-4}{}}
\newlabel{lemma-finite-pushforward-commutes-with-base-change@cref}{{[originallemma][0][]0959}{[1][0][]source}{}{}{}}
\makeatother
\title{Base change for the pushforward along an integral morphism}
\author{The Stacks Project Authors}
\date{stacks-project:etale-cohomology.tex:local-snapshot-20261011; GNU FDL 1.2 or later}
\begin{document}
\maketitle
\noindent\textit{Editorial note.} All delivered mathematical text is the author's own, copied from the pinned Stacks Project chapter file etale-cohomology.tex; no statement or proof step was rewritten, reordered or removed.\par\medskip
\noindent\textit{Editorial note.} The proof is the author's complete argument, including the reduction to the affine case and the descent of the local computation.\par\medskip
\noindent\textit{Editorial note.} Every labelled reference occurring inside the counted spans is carried in assembly.external\_references as a non-load-bearing navigation pointer (same-chapter ones at their definition site, cross-chapter ones at the invocation inside the counted span), so no load-bearing mathematics is replaced by a stub. No bibliography citation occurs inside the counted spans.\par\medskip
\noindent\textit{Editorial note.} The checked chapter list for same-name collisions is the set of stacks chapters staged in this repository: chow, cohomology, divisors, etale-cohomology, flat, more-algebra, more-morphisms, morphisms, spaces-more-morphisms, varieties together with this card's chapter etale-cohomology.\par\medskip
\noindent\textit{Editorial note.} The preamble is the chapter's own native preamble excerpt already used for this source: the Xy-pic setup (native preamble lines 14-18) and the native theorem declarations with the mathematical macros (native preamble lines 154-217).\par\medskip
\noindent\textit{Allowed standard prerequisites (editorial).}\begin{itemize}
\item Integral morphism of schemes and its properties. Same chapter together with the chapter Morphisms, the definition of an integral morphism and its stability under base change. Supplies the hypothesis on f and the affine nature of the morphism used in the reduction.
\item Pushforward and inverse image of sheaves of sets on the small etale site. Same chapter, the functors f\_*, g\textasciicircum{}\{-1\} and their behaviour on the small etale site. Supplies the two sides of the asserted equality and the comparison map.
\item Description of the small etale site of an integral morphism and its sections. Same chapter, the description of the pushforward along an integral morphism by the sections over the base change, used in the affine case. Supplies the computation at the end of the proof.
\end{itemize}\medskip
\ifdefined\hypertarget\hypertarget{sourcepointer-1}{}\fi
\noindent\textit{Source reference.} Original lemma 09YZ: sources/stacks-etale-cohomology/pinned-original-source.tex, source lines 8367--8367. Navigation pointer invoked at proof line 8367; the counted span contains the pointer only and the cited result is not reproduced inside it, so the pointer is registered as a non-load-bearing navigation stub in the established form for Stacks material. No bibliography citation is involved.\par\medskip
\ifdefined\hypertarget\hypertarget{sourcepointer-2}{}\fi
\noindent\textit{Source reference.} Original lemma 0EYN: sources/stacks-etale-cohomology/pinned-original-source.tex, source lines 7722--7722. Navigation pointer invoked at proof line 8371; the counted span contains the pointer only and the cited result is not reproduced inside it, so the pointer is registered as a non-load-bearing navigation stub in the established form for Stacks material. No bibliography citation is involved.\par\medskip
\ifdefined\hypertarget\hypertarget{sourcepointer-3}{}\fi
\noindent\textit{Source reference.} Original lemma 09Z1: sources/stacks-etale-cohomology/pinned-original-source.tex, source lines 7535--7535. Navigation pointer invoked at proof line 8375; the counted span contains the pointer only and the cited result is not reproduced inside it, so the pointer is registered as a non-load-bearing navigation stub in the established form for Stacks material. No bibliography citation is involved.\par\medskip
\ifdefined\hypertarget\hypertarget{sourcepointer-4}{}\fi
\noindent\textit{Source reference.} Original lemma 0959: sources/stacks-etale-cohomology/pinned-original-source.tex, source lines 8292--8292. Navigation pointer invoked at proof line 8399; the counted span contains the pointer only and the cited result is not reproduced inside it, so the pointer is registered as a non-load-bearing navigation stub in the established form for Stacks material. No bibliography citation is involved.\par\medskip
\begin{lemma}
\label{lemma-integral-pushforward-commutes-with-base-change}
Consider a cartesian square
$$
\xymatrix{
X' \ar[r]_{g'} \ar[d]_{f'} & X \ar[d]^f \\
Y' \ar[r]^g & Y
}
$$
of schemes with $f$ an integral morphism.
For any sheaf of sets $\mathcal{F}$ on $X_\etale$ we have
$f'_*(g')^{-1}\mathcal{F} = g^{-1}f_*\mathcal{F}$.
\end{lemma}

\begin{proof}
The question is local on $Y$ and hence we may assume $Y$ is affine.
Then we can write $X = \lim X_i$ with $f_i : X_i \to Y$ finite
(this is easy in the affine case, but see
Limits, Lemma \ref{limits-lemma-integral-limit-finite-and-finite-presentation}
for a reference). Denote $p_{i'i} : X_{i'} \to X_i$
the transition morphisms and $p_i : X \to X_i$ the projections.
Setting $\mathcal{F}_i = p_{i, *}\mathcal{F}$ we obtain
from Lemma \ref{lemma-linus-hamann}
a system $(\mathcal{F}_i, \varphi_{i'i})$
with $\mathcal{F} = \colim p_i^{-1}\mathcal{F}_i$.
We get $f_*\mathcal{F} = \colim f_{i, *}\mathcal{F}_i$
from Lemma \ref{lemma-relative-colimit}.
Set $X'_i = Y' \times_Y X_i$ with projections $f'_i$ and $g'_i$.
Then $X' = \lim X'_i$ as limits commute with limits.
Denote $p'_i : X' \to X'_i$ the projections. We have
\begin{align*}
g^{-1}f_*\mathcal{F}
& =
g^{-1} \colim f_{i, *}\mathcal{F}_i \\
& =
\colim g^{-1}f_{i, *}\mathcal{F}_i \\
& =
\colim f'_{i, *}(g'_i)^{-1}\mathcal{F}_i \\
& =
f'_*(\colim (p'_i)^{-1}(g'_i)^{-1}\mathcal{F}_i) \\
& =
f'_*(\colim (g')^{-1}p_i^{-1}\mathcal{F}_i) \\
& =
f'_*(g')^{-1} \colim p_i^{-1}\mathcal{F}_i \\
& =
f'_*(g')^{-1}\mathcal{F}
\end{align*}
as desired. For the first equality see above.
For the second use that pullback commutes with colimits.
For the third use the finite case, see
Lemma \ref{lemma-finite-pushforward-commutes-with-base-change}.
For the fourth use Lemma \ref{lemma-relative-colimit}.
For the fifth use that $g'_i \circ p'_i = p_i \circ g'$.
For the sixth use that pullback commutes with colimits.
For the seventh use $\mathcal{F} = \colim p_i^{-1}\mathcal{F}_i$.
\end{proof}

\end{document}
